%\documentclass[handout]{beamer}
\documentclass{beamer}\usepackage[]{graphicx}\usepackage[]{color}
%% maxwidth is the original width if it is less than linewidth
%% otherwise use linewidth (to make sure the graphics do not exceed the margin)
\makeatletter
\def\maxwidth{ %
  \ifdim\Gin@nat@width>\linewidth
    \linewidth
  \else
    \Gin@nat@width
  \fi
}
\makeatother

\definecolor{fgcolor}{rgb}{0.345, 0.345, 0.345}
\newcommand{\hlnum}[1]{\textcolor[rgb]{0.686,0.059,0.569}{#1}}%
\newcommand{\hlstr}[1]{\textcolor[rgb]{0.192,0.494,0.8}{#1}}%
\newcommand{\hlcom}[1]{\textcolor[rgb]{0.678,0.584,0.686}{\textit{#1}}}%
\newcommand{\hlopt}[1]{\textcolor[rgb]{0,0,0}{#1}}%
\newcommand{\hlstd}[1]{\textcolor[rgb]{0.345,0.345,0.345}{#1}}%
\newcommand{\hlkwa}[1]{\textcolor[rgb]{0.161,0.373,0.58}{\textbf{#1}}}%
\newcommand{\hlkwb}[1]{\textcolor[rgb]{0.69,0.353,0.396}{#1}}%
\newcommand{\hlkwc}[1]{\textcolor[rgb]{0.333,0.667,0.333}{#1}}%
\newcommand{\hlkwd}[1]{\textcolor[rgb]{0.737,0.353,0.396}{\textbf{#1}}}%
\let\hlipl\hlkwb

\usepackage{framed}
\makeatletter
\newenvironment{kframe}{%
 \def\at@end@of@kframe{}%
 \ifinner\ifhmode%
  \def\at@end@of@kframe{\end{minipage}}%
  \begin{minipage}{\columnwidth}%
 \fi\fi%
 \def\FrameCommand##1{\hskip\@totalleftmargin \hskip-\fboxsep
 \colorbox{shadecolor}{##1}\hskip-\fboxsep
     % There is no \\@totalrightmargin, so:
     \hskip-\linewidth \hskip-\@totalleftmargin \hskip\columnwidth}%
 \MakeFramed {\advance\hsize-\width
   \@totalleftmargin\z@ \linewidth\hsize
   \@setminipage}}%
 {\par\unskip\endMakeFramed%
 \at@end@of@kframe}
\makeatother

\definecolor{shadecolor}{rgb}{.97, .97, .97}
\definecolor{messagecolor}{rgb}{0, 0, 0}
\definecolor{warningcolor}{rgb}{1, 0, 1}
\definecolor{errorcolor}{rgb}{1, 0, 0}
\newenvironment{knitrout}{}{} % an empty environment to be redefined in TeX

\usepackage{alltt}
% ***************************************************************
% for handout, change only this...
%   \documentclass[twocolumn]{article}
%   \usepackage{beamerarticle}
%   \setlength{\textwidth}{7.5in}
%   \setlength{\textheight}{9.8in}
%   \setlength{\topmargin}{-1in}
%   \setlength{\oddsidemargin}{-.52in}
%   \setlength{\evensidemargin}{-.52in}

%\usepackage{beamerprosper}
%\usetheme{Warsaw}
%\usecolortheme{orchid}

\usepackage{graphicx}

\usepackage{amsmath,amssymb,array,comment,eucal}
\usepackage{amsmath,amssymb,array,eucal}
\usepackage{xcolor}
\definecolor{beamer@blendedblue}{RGB}{86,155,189}
\definecolor{myblue}{RGB}{12,76,138}
\setbeamercolor{structure}{fg=myblue}
\definecolor{Ftitle}{RGB}{12,76,138}
\definecolor{Descitem}{RGB}{238,238,244}
\definecolor{StdTitle}{RGB}{12,76,138}
\definecolor{StdBody}{RGB}{213,24,0}
\definecolor{StdBody}{RGB}{213,24,0}

\definecolor{AlTitle}{RGB}{255, 190, 190}
\definecolor{AlBody}{RGB}{213,24,0}

\definecolor{ExTitle}{RGB}{201, 217, 217}
\definecolor{ExBody}{RGB}{213,24,0}

\setbeamercolor{frametitle}{fg = Ftitle}
\setbeamercolor{title}{fg = Ftitle}
\setbeamercolor{item}{fg = Ftitle}
\setbeamercolor{subitem}{fg = Ftitle}
\setbeamercolor{subsubitem}{fg = Ftitle}
\setbeamercolor{description item}{fg = myblue}
\setbeamercolor{titlelike}{fg=myblue}
\newcommand{\e}{\mathbf{e}}
\renewcommand{\P}{\textsf{P}}
\newcommand{\R}{\textsf{R}}
\newcommand{\mat}[1] {\mathbf{#1}}
%\newcommand{\ind}{\mathrel{\mathop{\sim}\limits^{\mathit{ind}}}}
%\newcommand{\iid}{\mathrel{\mathop{\sim}\limits^{\mathit{iid}}}}
\newcommand{\E}{\textsf{E}}
\newcommand{\SE}{\textsf{SE}}
\newcommand{\SSE}{\textsf{SSE}}
\renewcommand{\SS}{\textsf{SS}}
\newcommand{\MSE}{\textsf{MSE}}
\newcommand{\SSR}{\textsf{SSR}}
\newcommand{\Be}{\textsf{Beta}}
\newcommand{\St}{\textsf{St}}
%\newcommand{\C}{\textsf{C}}
\newcommand{\GDP}{\textsf{GDP}}
\newcommand{\NcSt}{\textsf{NcSt}}
\newcommand{\Bin}{\textsf{Bin}}
\newcommand{\NB}{\textsf{NegBin}}
\renewcommand{\NG}{\textsf{NG}}
\newcommand{\N}{\textsf{N}}
\newcommand{\Ber}{\textsf{Ber}}
\newcommand{\Poi}{\text{Poi}}
\newcommand{\Gam}{\textsf{Gamma}}
\newcommand{\Gm}{\textsf{G}}
\newcommand{\Un}{\textsf{Unif}}
\newcommand{\Ex}{\textsf{Exp}}
\newcommand{\DE}{\textsf{DE}}
\newcommand{\tr}{\textsf{tr}}
\newcommand{\cF}{{\cal{F}}}
\newcommand{\cL}{{\cal{L}}}
\newcommand{\cI}{{\cal{I}}}
\newcommand{\cB}{{\cal{B}}}
\newcommand{\cP}{{\cal{P}}}
\newcommand{\bbR}{\mathbb{R}}
\newcommand{\bbN}{\mathbb{N}}
\newcommand{\pperp}{\mathrel{{\rlap{$\,\perp$}\perp\,\,}}}
\newcommand{\OFP}{(\Omega,\cF, \P)}
\newcommand{\eps}{\boldsymbol{\epsilon}}
\newcommand{\1}{\mathbf{1}_n}
\newcommand{\gap}{\vspace{8mm}}
\newcommand{\ind}{\mathrel{\mathop{\sim}\limits^{\rm ind}}}
\newcommand{\simiid}{\ensuremath{\mathrel{\mathop{\sim}\limits^{\rm
iid}}}}
\newcommand{\eqindis}{\ensuremath{\mathrel{\mathop{=}\limits^{\rm D}}}}
\newcommand{\iid}{\textit{i.i.d.}}
\newcommand{\SSZ}{S_{zz}}
\newcommand{\SZW}{S_{zw}}
\newcommand{\Bias}{\textsf{Bias}}
\newcommand{\Var}{\textsf{Var}}
\newcommand{\corr}{\textsf{corr}}
\newcommand{\diag}{\textsf{diag}}
\newcommand{\var}{\textsf{var}}
\newcommand{\Cov}{\textsf{Cov}}
\newcommand{\Sam}{{\cal S}}
\def\H{\mathbf{H}}
\newcommand{\I}{\mathbf{I}}
\newcommand{\Y}{\mathbf{Y}}
\newcommand{\tY}{\tilde{\mathbf{Y}}}
\newcommand{\Yhat}{\hat{\mathbf{Y}}}
\newcommand{\Yobs}{\mathbf{Y}_{{\cal S}}}
\newcommand{\barYobs}{\bar{Y}_{{\cal S}}}
\newcommand{\barYmiss}{\bar{Y}_{{\cal S}^c}}
\def\bv{\mathbf{b}}
\def\X{\mathbf{X}}
\def\tX{\tilde{\mathbf{X}}}
\def\x{\mathbf{x}}
\def\xbar{\bar{\x}}
\def\Xg{\mathbf{X}_{\boldsymbol{\gamma}}}
\def\Ybar{\bar{Y}}
\def\ybar{\bar{y}}
\def\y{\mathbf{y}}
\def\Yf{\mathbf{Y_f}}
\def\W{\mathbf{W}}
\def\w{\mathbf{w}}
\def\U{\mathbf{U}}
\def\V{\mathbf{V}}
\def\Q{\mathbf{Q}}
\def\Z{\mathbf{Z}}
\def\z{\mathbf{z}}
\def\v{\mathbf{v}}
\def\u{\mathbf{u}}

\def\zero{\mathbf{0}}
\def\one{\mathbf{1}}
\newcommand{\taub}{\boldsymbol{\tau}}
\newcommand{\betav}{\boldsymbol{\beta}}
\newcommand{\alphav}{\boldsymbol{\alpha}}
\newcommand{\A}{\mathbf{A}}
\def\a{\mathbf{a}}
\newcommand{\B}{\mathbf{B}}
\def\b{\boldsymbol{\beta}}
\def\bhat{\hat{\boldsymbol{\beta}}}
\def\tb{\tilde{\boldsymbol{\beta}}}
\def\bg{\boldsymbol{\beta_\gamma}}
\def\bgnot{\boldsymbol{\beta_{(-\gamma)}}}
\def\mub{\boldsymbol{\mu}}
\def\tmub{\tilde{\boldsymbol{\mu}}}
\def\muhat{\hat{\boldsymbol{\mu}}}
\def\t{\boldsymbol{\theta}}
\def\tk{\boldsymbol{\theta}_k}
\def\tj{\boldsymbol{\theta}_j}
\def\Mk{\boldsymbol{{\cal M}}_k}
\def\M{{{\cal M}}}
\def\Mj{{{\cal M}}_j}
\def\Mi{{{\cal M}}_i}
\def\Mg{{\cal M}_\gamma}
\def\Mnull{{\cal M}_{N}}
\def\gMPM{\boldsymbol{\gamma}_{\text{MPM}}}
\def\gHPM{\boldsymbol{\gamma}_{\text{HPM}}}
\def\Mfull{\boldsymbol{{\cal M}}_{F}}
\def\tg{\boldsymbol{\theta}_{\boldsymbol{\gamma}}}
\def\g{\boldsymbol{\gamma}}
\def\eg{\boldsymbol{\eta}_{\boldsymbol{\gamma}}}
\def\G{\mathbf{G}}
\def\cM{\cal M}
\def\D{\Delta}
\def \shat{{\hat{\sigma}}^2}
\def\uv{\mathbf{u}}
\def\l {\lambda}
\def\d{\delta}
\def\Sigmab{\boldsymbol{\Sigma}}
\def\Lambdab{\boldsymbol{\Lambda}}
\def\lambdab{\boldsymbol{\lambda}}
\def\Mg{{\cal M}_\gamma}
\def\S{{\cal{S}}}
\def\qg{p_{\boldsymbol{\gamma}}}
\def\pg{p_{\boldsymbol{\gamma}}}
\def\t{\boldsymbol{\theta}}
\def\T{\boldsymbol{\Theta}}

\title{Bayesian Model Averaging}

\author{Hoff Chapter 9, Hoeting et al 1999, Clyde \& George 2004, Liang et al 2008}
\date{\today}
%\Logo(-1.9,7.3){\includegraphics[width=.5in]{../eps/duke}}
% Optional: text to put in the bottom of each slide.
% By default, the title of the talk will be placed there.
%\slideCaption{\textit{October 28, 2005 }}
\IfFileExists{upquote.sty}{\usepackage{upquote}}{}
\begin{document}
% make the title slide
\maketitle

\begin{frame} \frametitle{Outline}

\begin{itemize}
\item Problems with g-priors
\item Alternatives:  Mixtures of $g$-priors
\item Model Averaging
\item Choice of Model
\end{itemize}

\end{frame}

\begin{frame}
\frametitle{Bayes Factors}
\begin{itemize}
\item Bayes Factor = ratio of marginal likelihoods \pause

\item Posterior odds = Bayes Factor $\times$ Prior odds \pause

\item Posterior Probability
$$P(\Mg \mid \Y) = \frac{BF[\Mg : \M_0] p(\Mg)/p(\M_0)}
{\sum_{\Mg \in \Gamma} BF[\Mg : \M_0] p(\Mg)/p(\M_0)}
$$

\end{itemize}
\end{frame}

\begin{frame}  \frametitle{Problem with g-Prior with arbitrary g}

The Bayes factor for comparing $\Mg$ to the null
model:
%$$
% BF(\Mg : \M_0) =    (1 + g)^{(n - 1 - p_\g)/2} (1 + g(1 - \R^2))^{-(n-1)/2}
%$$
\pause
\begin{itemize}
\item Let $g$ be a fixed constant and take $n$ fixed.\pause
\item Let $F = \frac{R_{\g}^2/\pg}{(1 - R_{\g}^2)/(n - 1 - \pg)}$ \pause
\item As $R^2_{\g} \to 1$, $F \to \infty$ LR test would reject $H_0$  where
 $F$ is the usual $F$ statistic for  comparing model $\Mg$ to $\M_0$ \pause
\item Bayes Factor would go to $(1 + g)^{(n-\pg -1)/2}$ as $F \to \infty$  (bounded for fixed $g$, $n$ and $\pg$ \pause
\end{itemize}

Bayes and Frequentist would not agree in this limit \pause

\vspace{12pt}
``Information paradox''
\end{frame}


\begin{frame} \frametitle{Resolution of Paradox}

Liang et al (2008) show that paradox can be resolved with mixtures of $g$-priors
$$p(\bg \mid \phi) = \int_0^{\infty} \N(\bg; 0, g (\Xg^T\Xg)^{-1}/\phi) p(g) dg$$
\pause

\begin{itemize}
\item  $BF \to \infty$ if $\R^2 \to 1$  $\Leftrightarrow$ $\E_g[(1 +
g)^{-\pg/2}]$ diverges \pause
\item Zellner-Siow Cauchy prior
$$ 1/g \sim \Gam(1/2, n/2)$$ \pause

\item Hyper-g   $p(g) \propto (1 +g )^{a/2 - 1}$ if $2 < a \le 3$

$$\frac{g}{1+g} \sim \textsf{Beta}(1, \frac{a}{2} - 1)$$
\pause
\item "hyper-g/n"
\pause
\item  robust prior (Bayarri et al Annals of Statistics 2012)
\end{itemize}

\end{frame}


\begin{frame}[fragile] \frametitle{Example}


\begin{knitrout}
\definecolor{shadecolor}{rgb}{0.969, 0.969, 0.969}\color{fgcolor}\begin{kframe}
\begin{alltt}
\hlkwd{library}\hlstd{(BAS)}
\hlstd{poll.ZS} \hlkwb{=} \hlkwd{bas.lm}\hlstd{(}\hlkwd{log}\hlstd{(SO2)} \hlopt{~} \hlstd{temp} \hlopt{+} \hlkwd{log}\hlstd{(firms)} \hlopt{+}
                             \hlkwd{log}\hlstd{(popn)} \hlopt{+} \hlstd{wind} \hlopt{+}
                             \hlstd{precip}\hlopt{+} \hlstd{rain,}
                  \hlkwc{data}\hlstd{=usair,}
                  \hlkwc{prior}\hlstd{=}\hlstr{"ZS-null"}\hlstd{,}
                  \hlkwc{alpha}\hlstd{=}\hlnum{41}\hlstd{,}    \hlcom{# g = n}
                  \hlkwc{n.models}\hlstd{=}\hlnum{2}\hlopt{^}\hlnum{7}\hlstd{,}\hlcom{# enumerate (can omit)}
                  \hlkwc{modelprior}\hlstd{=}\hlkwd{uniform}\hlstd{(),}
                  \hlkwc{method}\hlstd{=}\hlstr{"deterministic"}\hlstd{)}   \hlcom{# fast enumeration}
\end{alltt}
\end{kframe}
\end{knitrout}
 use `prior = "hyper-g"` and `a = 3` for hyper-g or `prior = "hyper-g/n"` and `a=3` for hyper-g/n
\end{frame}

\begin{frame}
\begin{knitrout}
\definecolor{shadecolor}{rgb}{0.969, 0.969, 0.969}\color{fgcolor}\begin{kframe}
\begin{alltt}
\hlkwd{plot}\hlstd{(poll.ZS,} \hlkwc{which}\hlstd{=}\hlnum{4}\hlstd{)}
\end{alltt}
\end{kframe}

{\centering \includegraphics[width=\maxwidth]{figure/unnamed-chunk-2-1} 

}



\end{knitrout}
\end{frame}

\begin{frame} \frametitle{Bayesian Model Averaging}

\begin{itemize}
\item Posterior for $\mu = \one \alpha + \X \b$ is a mixture distribution

$$p(\mu \mid \Y) = \sum p(\mu \mid \Y, \Mg) p(\Mg \mid \Y)$$ \pause
with expectation expressed as a weighted average
$$\E[\mu \mid \Y] = \one \hat{\alpha} +  \X \sum \E[\b \mid \Y, \Mg] p(\Mg \mid \Y)$$
\pause
\item Predictive Distribution for $\Y^*$
$$p(\Y^* \mid \Y) = \sum p(\Y^* \mid \Y, \Mg) p(\Mg \mid \Y)$$
\pause
\item Posterior Distribution of $\beta_j$

$$p(\beta_j \mid \Y) = p(\gamma_j = 0 \mid \Y) \delta_0(\beta) + \sum p(\beta_j \mid \Y, \Mg) \gamma_j p(\Mg \mid \Y) $$
\end{itemize}

\end{frame}

\begin{frame} \frametitle{Estimator}

\begin{itemize}
\item Find $\hat{\mu}$ that minimizes posterior expected loss
$$\E[(\mu - \hat{\mu})^T(\mu - \hat{\mu}) \mid \Y]$$ \pause
\item Solution is posterior mean under BMA
$$\E[\mu \mid \Y] = \one \hat{\alpha} +  \X \sum \E[\b \mid \Y, \Mg] p(\Mg \mid \Y)$$ \pause
\item If one model has probability 1, then BMA is equivalent to using the highest posterior probability model \pause
\item incorporates estimates from other models when there is substantial uncertainty
\end{itemize}

\end{frame}

\begin{frame}[fragile] {Coefficients under BMA}
\begin{knitrout}
\definecolor{shadecolor}{rgb}{0.969, 0.969, 0.969}\color{fgcolor}\begin{kframe}
\begin{alltt}
\hlstd{beta.ZS} \hlkwb{=} \hlkwd{coef}\hlstd{(poll.ZS)}
\hlstd{beta.ZS}
\end{alltt}
\begin{verbatim}
## 
##  Marginal Posterior Summaries of Coefficients: 
## 
##  Based on the top  64 models 
##             post mean  post SD    post p(B != 0)
## Intercept    3.153004   0.082226   1.000000     
## temp        -0.058053   0.020325   0.976833     
## log(firms)   0.206384   0.177253   0.758554     
## log(popn)   -0.035074   0.174760   0.342677     
## wind        -0.129875   0.085195   0.813330     
## precip       0.010898   0.011327   0.633639     
## rain         0.001759   0.004034   0.356085
\end{verbatim}
\end{kframe}
\end{knitrout}
\end{frame}

\begin{frame}[fragile] \frametitle{Posterior of Coefficients under BMA}
\begin{knitrout}
\definecolor{shadecolor}{rgb}{0.969, 0.969, 0.969}\color{fgcolor}

{\centering \includegraphics[width=\maxwidth]{figure/coefplot-1} 

}



\end{knitrout}
\end{frame}

\begin{frame}[fragile] \frametitle{Credible Intervals for Coefficients under BMA}
\begin{knitrout}
\definecolor{shadecolor}{rgb}{0.969, 0.969, 0.969}\color{fgcolor}\begin{kframe}
\begin{alltt}
\hlkwd{plot}\hlstd{(}\hlkwd{confint}\hlstd{(beta.ZS,} \hlkwc{parm}\hlstd{=}\hlnum{2}\hlopt{:}\hlnum{7}\hlstd{))}
\end{alltt}
\end{kframe}

{\centering \includegraphics[width=\maxwidth]{figure/unnamed-chunk-4-1} 

}


\begin{kframe}\begin{verbatim}
## NULL
\end{verbatim}
\end{kframe}
\end{knitrout}
\end{frame}

\begin{frame} \frametitle{Selection and Model Uncertainty}

\begin{itemize}
\item Select a model and $\hat{\mu}$ that minimizes posterior expected loss
$$\E[(\mu - \hat{\mu})^T(\mu - \hat{\mu}) \mid \Y]$$ \pause
\item BMA is "best" estimator without selection \pause

\item Best model and estimator is the posterior mean under the model that is closest to BMA under squared error loss
$$(\hat{\mu}_{BMA} - \hat{\mu}_{\Mg})^T(\hat{\mu}_{BMA} - \hat{\mu}_{\Mg})$$
\pause
\item Often contains more predictors than the HPM or Median Probability Model

\end{itemize}
\end{frame}

\begin{frame}[fragile] \frametitle{Best Predictive Model}
\begin{knitrout}
\definecolor{shadecolor}{rgb}{0.969, 0.969, 0.969}\color{fgcolor}\begin{kframe}
\begin{alltt}
\hlcom{#BPM}
\hlstd{BPM} \hlkwb{=} \hlkwd{predict}\hlstd{(poll.ZS,} \hlkwc{estimator} \hlstd{=} \hlstr{"BPM"}\hlstd{)}
\hlstd{BPM}\hlopt{$}\hlstd{bestmodel}
\end{alltt}
\begin{verbatim}
## [1] 0 1 2 4 5 6
\end{verbatim}
\begin{alltt}
\hlstd{(poll.ZS}\hlopt{$}\hlstd{namesx[}\hlkwd{attr}\hlstd{(BPM}\hlopt{$}\hlstd{fit,} \hlstr{'model'}\hlstd{)} \hlopt{+}\hlnum{1}\hlstd{])[}\hlopt{-}\hlnum{1}\hlstd{]}
\end{alltt}
\begin{verbatim}
## [1] "temp"       "log(firms)" "wind"       "precip"     "rain"
\end{verbatim}
\begin{alltt}
\hlcom{#HPM}
\hlstd{HPM} \hlkwb{=} \hlkwd{predict}\hlstd{(poll.ZS,} \hlkwc{estimator} \hlstd{=} \hlstr{"HPM"}\hlstd{)}
\hlstd{HPM}\hlopt{$}\hlstd{bestmodel}
\end{alltt}
\begin{verbatim}
## [1] 0 1 2 4 5
\end{verbatim}
\end{kframe}
\end{knitrout}


\end{frame}



\begin{frame}[fragile] %\frametitle{Comparison}
\begin{knitrout}
\definecolor{shadecolor}{rgb}{0.969, 0.969, 0.969}\color{fgcolor}

{\centering \includegraphics[width=\maxwidth]{figure/plot_est-1} 

}



\end{knitrout}
\end{frame}

\begin{frame} \frametitle{Summary}

\begin{itemize}
\item BMA shown in practice to have better out of sample predictions than selection (in many cases) \pause
\item avoids selecting a single model and accounts for out uncertainty \pause
\item if one model dominates BMA is very close to selection  (asymptotically will put probability one on model that is "closest" to the true model) \pause
\item MCMC allows one to implement without enumerating all models \pause
\item BMA depends on prior on coefficients, variance and models  (sensitivity to choice?) \pause
\item Mixtures of $g$ priors preferred to usual $g$ prior
\end{itemize}
\end{frame}
\end{document}
